\chapter{}

Pietro era doventato così negligente, che verso il mese di maggio il rettore non lo volle più alla scuola.

Domenico lo percosse con lo scheggiale dei calzoni, fino a far piangere anche Anna. Ma, il giorno dopo, nessuno gli disse più niente.

Anna spiegò a Rebecca:

— Sono le imprecazioni di quelli che ci vogliono male.

Fece tutti i giorni alcune preghiere ad un santo; ma non trovò mai modo di parlarne sul serio al marito, che le rispondeva sempre:

— Oggi non posso.

Se lo tratteneva per la giubba, egli la lasciava con queste parole:

— Pensaci tu a lui. Anche tu ora....

Ella non osava di più, temendo che se la rifacesse con Pietro; stordendolo a forza di pugni, con il pretesto di essersi arrabbiato anche troppo.

Nè meno la notte era possibile, perchè a pena gliene discorreva, stringeva i pugni e gridava:

— Lasciami dormire. Ho sonno; è da stamani che lavoro. Riposati anche tu....

Oppure rispondeva:

— Hai contato bene i denari incassati oggi? Prima di venire a letto, dovevi contarli. È necessario.

Se ella, per rendergli il cambio, stava zitta, le alzava il capo dal guanciale:

— Rispondi!

Aspettava un poco, tentando di questionare; ma poi si addormentava.

Durante una loro contesa in bottega, Pietro saltò fuori a dire:

— Imparerò il disegno.

Lo scritturale di un notaio, che aveva finito allora di mangiare, fece una enorme risata.

Pietro lo guardò a lungo, sbigottito dei suoi occhi dolci e contenti che lo compativano.

Era un uomo grasso; dal volto lucido e purpureo, sparso di bitorzoli. Aveva un vestito chiaro e una catena d'oro; i capelli biondicci, la fronte bassa. Disse a Domenico, con convinzione tranquilla:

— Non gli date retta. Fategli imparare il vostro mestiere. Voi trattori guadagnate quanto volete.

Tutti risero, perchè alludeva al conto che doveva pagare.

Pietro, mentre una specie di formicolìo lieve attraversava il suo volto, dal mento alla fronte, esclamò:

— Che importa a lei?

Costui trasse da un astuccio di cuoio un bocchino d'ambra cerchiata d'oro, e v'infilò mezzo sigaro. Poi disse:

— Vai a comprarmi una scatola di fiammiferi.

E gli gettò un soldo su la tavola.

Pietro guardò anche suo padre: tutti lo fissavano; i volti e gli occhi bruciavano la sua anima. Il cuore gli batteva.

Domenico disse:

— Vai, dunque!

Egli afferrò la moneta, e corse dal tabaccaio.

Allora lo scritturale rise tanto che fece il viso congestionato; e, tra gli scoppi di tosse, aggiunse:

— Fatelo ubbidire più che potete.

Anna soffriva di queste domestichezze; ma, per paura di perdere gli avventori, non ci si metteva a tu per tu. Invece Domenico se n'esaltava; e gli pareva sempre più di aver ragione. E diceva a Pietro:

— Stai attento a quello che ti dico io. Non hai più bisogno di studiare. Basta che tu sappia fare la moltiplicazione. Dovrebbero esser abolite le scuole, e mandati tutti gli insegnanti a vangare. La terra è la migliore cosa che Dio ci ha data.

Anna, scontenta, rispondeva:

— Codeste sono idee tue.

Domenico chiedeva, con scherno:

— Quanto tempo ci sei andata a scuola tu?

Non ci mancava che da contrastare con la moglie! Ella scuoteva la testa.

— Noi, senza saper né meno la nostra firma, abbiamo fatto fortuna.

Gli avventori rimanevano pensosi; poi esclamavano, tanto per non scontentare di più Anna:

— È ancora giovine. Non c'è da capire quel che ci potrete ricavare.

— Ma anche quando io avrò sessant'anni, ed egli più di venti, sarò sempre capace di rompergli la testa.

— Oh, grosso e forte come voi non verrà

di certo!
